% ============================================================
% declaracion-ia.tex — Uso de IA
% SemArq1_TeamBales.tex, DESPUÉS de \printbibliography, como exige la
% orientación del seminario (sección 3): epígrafe independiente, al
% final del documento, que indique (1) herramienta y versión,
% (2) propósito y (3) forma de integración del resultado.
% Fuente de trabajo: «Uso de IA (Claude).md».
% ============================================================

\chapter*{Uso de IA}
\addcontentsline{toc}{chapter}{Uso de IA}

Para la elaboración de este seminario, el equipo utilizó dos herramientas de inteligencia artificial. Heily Rodríguez Rodríguez y Ernesto Alejandro Soler Choong usaron \textbf{Claude}, de Anthropic, con el modelo \textbf{Claude Opus 5.5}. Enrique González Moreira usó \textbf{OpenCode} con el modelo \textbf{GLM 5.3}.

Las herramientas se emplearon para buscar en la bibliografía recomendada los capítulos, epígrafes y páginas útiles para cada pregunta, y para comprobar que cada cita remite a la página o sección correcta. También se usaron para redactar borradores de las respuestas, mejorar la formalidad y la claridad de la redacción, y traducir al español ideas y fragmentos de las fuentes, que en su mayoría están en inglés. Por último, sirvieron para ayudar a cumplir los requisitos formales de la orientación, como la portada, el estilo de citas y la correspondencia entre las citas del texto y las referencias, y para generar el código \LaTeX{} del documento, incluidos los diagramas.

Los borradores generados se utilizaron como base de las respuestas, que el equipo revisó, corrigió y reescribió hasta su versión final. Las citas se compararon con los textos originales de la bibliografía. Las traducciones sirvieron para comprender y parafrasear las ideas de los autores, y los fragmentos traducidos que aparecen entre comillas se contrastaron con el texto original. El código \LaTeX{} y los diagramas se revisaron y ajustaron hasta obtener el resultado deseado.

Todo el contenido generado fue revisado, verificado y adaptado por el equipo, quien asume la responsabilidad final del documento.
